% Page budget: 1.55 pages | Owns: negation, OBC derivation, additional-state scope, implementation, and true-parameter condition.
% WRITING GUIDE
% Purpose: Explain negation, derive the final state-level orthogonal-by-construction map, and delimit its interpretability guarantee.
% Sources: README "Source map per subsection", rows IV-A, IV-B and IV-C; mandatory papers and the reference gate as in the README.
% Research-plan anchor: Preserve Aspect 3's goal of protecting physical interpretability; state clearly how the final construction differs from the proposed regularisation.
% Current decisions: Reduced identifiable coordinates, data-derived reference tuples, a tangent-only one-step physical-state projection, per-epoch refresh, and no affine comparison arm.
% Choices to justify: Protected parameter coordinates, reference-set construction, projected rows, rank tolerance, exclusion of the affine offset, refresh frequency, and treatment of additional states.
% Cautions: docs/orthogonal-projection-plan.md describes an older soft-penalty route; one-step state orthogonality is not trajectory-output orthogonality or guaranteed true recovery.
% Planned figures: New geometric projection figure with a physical and additional-state partition inset; show the remaining uniqueness limitation as clearly as the protected direction.
% Section is finished when: The protected space, coordinate frame, gradient treatment, added-state scope, and true-recovery condition are explicit.
%
% DRAFT NOTES (cycle 14, 2026-10-09)
% MUST ESTABLISH (PROPOSED, inferred from the README ownership row IV and source-map rows IV-A to IV-D;
%   the file held no agreed block)
% Contribution delivered: the third contribution of Section I in its wording (state-level orthogonality for the
%   identifiable physical-parameter tangent space and its true-parameter recovery condition).
% IV-A 1. Negation (Kessels Sec. 5.3.1.3) made precise for the parameter directions: to first order a learned write
%        -Phi dxi cancels a displacement dxi; Phi taken in the ten training coordinates xi (D-190 Why: raw
%        parameters give nullity 4). On the gantry these directions are forces scaled by M(Y)^{-1}, so they move with Y.
% IV-B 2. Data-derived reference tuples (P^{-T} q_act, fourth-order difference, xbar = 0, recorded force; D-192).
%      3. Gyorok 2026 Eqs. 8, 10, 11 adapted: the regressor is the Jacobian of the nonlinear RK4 transition at xi°
%        (Gyorok 2025 Sec. 4 precedent), tangent only (D-186 of 2026-09-10 reasons 1, 2 and its caveat);
%        coefficient differentiated through (D-190 item 5). Differs from the proposed penalty (research plan).
% IV-C 4. M_OBC: the problem display with the corrected physical line; xi° refreshed per epoch (D-190 item 6;
%        Gyorok 2025 Remark 2 adapted); per phase what changes (PS2 S1, S2, S3, polish, selection).
% IV-D 5. Proposition 1 (own): on the tuples and to first order, the estimate is xi* + Phi^+ eps*, unique for every
%        corrected write; recovery iff Phi' eps* = 0 (state-level Gyorok Cond. 4, Eq. 21). Scope paragraph.
% CORRECTIONS against the earlier draft of this file (code and sources win):
%   - The Jacobian is taken in the training coordinates xi, not in theta_base (FS/obc.py build_stacked_sensitivity
%     of transition_from_free); the range is the same because d theta_base / d xi is invertible (checked: diagonal
%     entries theta_j > 0 and b s sech^2 != 0, triangular otherwise, s != 0 needs m_Delta,0 != 0).
%   - Gyorok 2025 Sec. 4 does NOT say a fixed expansion point loses accuracy; it argues the nominal point stays
%     valid. The draft states only what the paper says.
%   - The recovery statement concerns the tuples whose successor has a full stencil (D-203 drops the last tuple of
%     each record), so it is named a variant on I'.
%   - Gyorok 2026 Eq. (39) does not support parameter uniqueness (citation-log row 74); Assumption 1, Eq. (7) used.
%   - Symbols: bar marks the added state and tilde the physical state, so the expansion point is xi° and the
%     corrected write f^perp; Delta names the scheduling block, so the true discrepancy is eps*; z names the LFR
%     latent variable, so the tuples are written by their state and input; + marks the pseudoinverse, so the
%     stacked successor is chi.
% LEFT OUT (reader test): hand-written forward sensitivity (D-194), per-objective hoist (D-195), float64 solve and
%   SVD storage, rank tolerance and the rank-loss fallback (mechanics), validation diagnostics (D-201), the affine
%   eleven-column arm (not in the campaign, THESIS-RESULTS), the overlap measures (THESIS-RESULTS reports none;
%   todo in IV-B), the stencil's gain argument in the old 140 to 230 Hz band (superseded data), the Section II
%   parameter prior (Section III's block: never mention it).
% EXISTING TODOS of the restored file: "Source" (opening) resolved by the contribution sentence and Gyorok 2026
%   Sec. 6; "Subsection title improve" resolved; the Section III negation todo resolved (IV-A); the affine-arm todo
%   and its notes resolved (no affine arm, THESIS-RESULTS; exclusion reasoned in IV-B); "one-step versus
%   trajectory penalty" resolved by D-190 Ruled out (1) (trajectory method rejected by the user); "orthogonal
%   multisine design" resolved by THESIS-RESULTS Sec. 6 (T_OA not run) and stated with Remark 5 in IV-D; the
%   Gamma todo kept and rewritten (supervisor confirmation, IV-B). Notes moved here: figure plan (geometric
%   decomposition of f_aug into its part in range(Phi) and the orthogonal part, inset with the state partition);
%   pointwise no-go and Coulomb inner-product argument belong to Section VII; rho* = 0.205 and the 4.96 percent
%   floor dropped (earlier absorber-only data).
% STALE LINES ELSEWHERE (not edited): 03_augmentation.tex line 667 claims Section IV makes the estimate unique,
%   while Proposition 1 gives uniqueness on the tuples and to first order only; 05_setup.tex line 52 lists an
%   affine-OBC arm; 05_setup.tex line 59 reports the Jacobian rank with respect to varphi (now xi);
%   06_results.tex line 40 cites eq:obc_overlap, removed here; 06_results.tex lines 45, 49, 50, 52 quote the
%   affine arm, rho* and the 4.96 percent floor of earlier data. 99_appendix.tex OBC bullets updated.
% CODE CONFIRMED this session: GD/obc_gantry.py (build_reference_set, GantryOBCCorrection.forward,
%   make_reference_field with xbar = 0, attach_obc), FS/obc.py (build_stacked_sensitivity, OBCBasis.coefficient,
%   OBCLifecycle.refresh, maybe_refresh, coefficient), FS/interconnect.py (forward subtraction, loss seam,
%   _sync_prediction_state_for_validation, rebuild after the selected checkpoint), FS/ps2_opening.py (S1 residual
%   and S2 through the prediction path with the corrected model), FS/lbfgs_polish.py (no OBC logic: xi° fixed
%   during the polish), entry _THESIS_FIXED (obc_space tangent, obc_refresh epoch).
% LENGTH: see the delivery report of cycle 14.
\section{Interpretability by Orthogonal Construction}\label{sec:obc}
\ifdraft
\begin{itemize}
\item Goal from Section~III-D: a unique estimate of $\theta_{\mathrm{base}}$; names IV-A to IV-D (style profile)
\item Contribution share: the third contribution of Section~I, in its wording (Introduction)
\end{itemize}
\fi

Section~\ref{sec:aug_closed_loop} requires a unique estimate of
$\theta_{\mathrm{base}}$, which the additive split of the augmented model does
not give.
Our goal is a correction of the learned physical write that removes the
parameter directions of the baseline from it.
To this end, we state the negation it removes
(Section~\ref{sec:obc_negation}), construct the
orthogonal-by-construction (OBC) correction (Section~\ref{sec:obc_map}) and
identify the corrected model (Section~\ref{sec:obc_id}).
We then derive when its estimate equals the true parameters and delimit the
guarantee (Section~\ref{sec:obc_scope}).
This section constructs state-level orthogonality for the identifiable
physical-parameter tangent space and quantifies its true-parameter recovery
condition, the third contribution of Section~I.

\subsection{Negation of the physical parameters}\label{sec:obc_negation}
\ifdraft
\begin{itemize}
\item Kessels' negation; the linear example of Gy\"or\"ok et al.\ (Kessels Sec.~5.3.1.3; Gy\"or\"ok 2025 Ex.~1)
\item Directions in the ten training coordinates $\xi$, since the raw parameters lose four; same range as in $\theta_{\mathrm{base}}$ (D-190; code; this work)
\item Sensitivity $\Phi_{\xi^\circ}$ displayed; to first order a learned write cancels a displacement (this work; first-order expansion)
\item On the gantry the directions are forces scaled by $M(Y)^{-1}$, so they change with $Y$ (this work; research plan Aspect~3)
\end{itemize}
\fi

The non-unique split of Section~\ref{sec:aug_closed_loop} needs a precise form
before a construction can remove it.
Kessels calls learned terms that (partially) negate the first-principles terms
negation~\cite[Sec.~5.3.1.3]{kessels2025ai}.
For a state transition linear in its parameter, Gy\"or\"ok et al.\ show that a
linear learned term makes every pair with the same sum a
minimiser~\cite[Ex.~1]{gyorok2025l4dc}.

These directions are the changes of the normalised transition
$f^{\mathrm n}_{\mathrm{base}}$ of \eqref{eq:aug_norm} that $\theta_{\mathrm{base}}$
can produce.
We take them in the training coordinates $\xi$ of \eqref{eq:mdelta_map}.
In the fourteen raw parameters, the four lost directions of
Section~\ref{sec:params} would make them linearly dependent.
Because $\partial\theta_{\mathrm{base}}/\partial\xi$ is invertible
(Appendix~\ref{app:lfr-wellposed}), they span the same directions as the
derivatives with respect to $\theta_{\mathrm{base}}$.
At an expansion point $\xi^\circ$, they are the columns of the sensitivity
\begin{equation}
 \Phi_{\xi^\circ}(\tilde x^{\mathrm n},u^{\mathrm n})
 =\left.\frac{\partial f^{\mathrm n}_{\mathrm{base}}\bigl(\theta_{\mathrm{base}}(\xi),
   \tilde x^{\mathrm n},u^{\mathrm n}\bigr)}{\partial\xi}\right|_{\xi=\xi^\circ},
 \label{eq:obc_sens}
\end{equation}
where $\xi^\circ\in\R^{10}$, $\Phi_{\xi^\circ}:\R^6\times\R^3\to\R^{6\times10}$,
and $\theta_{\mathrm{base}}(\xi)$ is the map \eqref{eq:mdelta_map}.
To first order, a displacement $\delta\xi$ changes the transition by
$\Phi_{\xi^\circ}\delta\xi$, which a learned write $-\Phi_{\xi^\circ}\delta\xi$
cancels, so the data cannot attribute this change to either component.

On the gantry, these directions are forces scaled by the inverse inertia.
Differentiating \eqref{eq:eom} at fixed $(q,\dot q,u)$ gives
$\partial\ddot q/\partial\theta_{\mathrm{base},j}
=-M(Y)^{-1}\bigl(\partial_jM(Y)\,\ddot q+\partial_jC\,\dot q+\partial_jK\,q\bigr)$,
with $\partial_j=\partial/\partial\theta_{\mathrm{base},j}$.
A write that negates a mass, damping or stiffness combination therefore acts as
an added inertial, damping or elastic force.
Through $M(Y)^{-1}$, the direction it negates changes with the payload
position.

\subsection{Orthogonal-by-construction correction}\label{sec:obc_map}
\ifdraft
\begin{itemize}
\item Removing negation needs the learned write on a fixed set of states; only positions are measured (Section~III-B)
\item Reference tuples from every training sample: $P^{-\top}\qact$, fourth-order central difference, $\bar x=0$, recorded force; stacks over all six physical rows (D-192; code); why not a baseline simulation (?)
\item Adapted from Gy\"or\"ok 2026: regressor is the Jacobian of the nonlinear RK4 step, as Gy\"or\"ok 2025 linearises (Gy\"or\"ok 2026 Eq.~(2); Gy\"or\"ok 2025 Sec.~4)
\item Tangent only: the offset charges no parameter direction and depends on the centring; nominal-response caveat (D-186 of 2026-09-10; D-192); supervisor confirmation (?)
\item Coefficient and corrected write displayed; no weight, unlike the proposed penalty (Gy\"or\"ok 2026 Eqs.~(8), (10), (11), Assumption~1; Gy\"or\"ok 2025 Remark~1; research plan Aspect~3)
\item Differentiated coefficient: the gradient of the corrected write stays orthogonal (D-190 item 5; this work, algebra); measured check (?)
\end{itemize}
\fi

Removing negation needs the part of the learned physical write in the range
of $\Phi_{\xi^\circ}$, measured over a fixed set of states.
Only the positions are measured (Section~\ref{sec:aug_encoder}), so the states
of this set must be reconstructed.
We build one reference tuple from every training sample at which a
fourth-order central difference fits inside its record.
The stacks run over all six physical rows, which the network writes
(Section~\ref{sec:aug_structure}):
{\small
\begin{equation}
\begin{aligned}
 \boldsymbol\Phi&=\col\bigl(\Phi_{\xi^\circ}(\tilde x^{\mathrm n}_i,
   u^{\mathrm n}_i)\bigr)_{i\in\mathcal I},\\
 \boldsymbol f_{\mathrm{aug}}(\theta_{\mathrm{aug}})&=\col\bigl(
   f^{\mathrm n}_{\mathrm{aug}}(\theta_{\mathrm{aug}},
   \col(\tilde x^{\mathrm n}_i,0),u^{\mathrm n}_i)\bigr)_{i\in\mathcal I},\\
 \tilde x^{\mathrm n}_i&=T_x\bigl(\col(q_i,\dot q_i)-\mu_x\bigr),\qquad
 q_i=P^{-\top}q_{\mathrm{act},i},
\end{aligned}
\label{eq:obc_stacks}
\end{equation}
}
where $\mathcal I$ is the set of these samples, $N=|\mathcal I|$,
$\boldsymbol\Phi\in\R^{6N\times10}$,
$\boldsymbol f_{\mathrm{aug}}:\R^{n_{\mathrm{aug}}}\to\R^{6N}$, $\dot q_i$ is
the fourth-order central difference of $q$ over $T_s$ at sample $i$, and
$u^{\mathrm n}_i$ is the normalised recorded force, with $T_x$ and $\mu_x$ of
\eqref{eq:aug_norm}.
The tuples depend neither on the current model nor on the controller, so they
stay fixed for the whole identification.
\todo{Missing: (a) reason for data-derived tuples rather than a simulation of the baseline on the training data, which Gy\"or\"ok et al.\ use when the states are not measured~\cite[Sec.~3]{gyorok2025l4dc}. Candidate: an open-loop simulation drifts on the free axes (Section~\ref{sec:aug_loop}); D-185 records this as the reason of D-111, both superseded entries, so it needs Dirk's confirmation. (b) On the noisy records, the reference basis built from filtered positions spans nearly the same range as the one from noise-free positions (D-232, OBC velocity result), but no section reports it. Candidate: one sentence with a pointer to the appendix, or leave out.}

We adapt the construction of Gy\"or\"ok et al., whose regressor belongs to an
input-output baseline linear in its parameters~\cite[Eq.~(2)]{gyorok2026obc}.
Our transition is rational in $\theta_{\mathrm{base}}$ through $M(Y)^{-1}$ and
has no such form.
We therefore use the sensitivity stack, as their penalty method linearises a
nonlinear baseline at an expansion point~\cite[Sec.~4]{gyorok2025l4dc}.
Unlike there, we keep only the tangent and omit the offset column of the
expansion~\cite[Eq.~(18)]{gyorok2025l4dc}.
To first order, a parameter change moves the transition only within
$\ran\boldsymbol\Phi$ (Section~\ref{sec:obc_negation}).
The offset column would therefore charge a direction that no parameter change
produces.
It also depends on the centring $\mu_x$ of \eqref{eq:aug_norm}, while
$\boldsymbol\Phi$ does not.
A write that cancels the nominal transition is consequently charged only
through its component in $\ran\boldsymbol\Phi$.
\todo{Missing: supervisor confirmation of the protected set. D-190 recorded the supervisors' condition as non-overlap with the baseline output, which the omitted offset protects only partly. Candidate: the tangent set, for the reasons above (D-192; D-186 of 2026-09-10).}

Gy\"or\"ok et al.\ remove from the stacked learned output its least-squares fit
on the regressor, whose coefficient is a function of the network
parameters~\cite[Eqs.~(8), (10), (11)]{gyorok2026obc}.
Unlike their earlier penalty, which the research plan proposed, this removes
the overlap without a weight that trades it against the
fit~\cite[Remark~1]{gyorok2025l4dc}.
With $\boldsymbol\Phi$ of full column rank, as their Assumption~1
requires~\cite[Eq.~(7)]{gyorok2026obc}, the construction is
\begin{equation}
\begin{aligned}
 \theta_{\mathrm{aux}}&=\boldsymbol\Phi^{+}\boldsymbol f_{\mathrm{aug}}(\theta_{\mathrm{aug}}),\\
 \boldsymbol f^{\perp}_{\mathrm{aug}}(\theta_{\mathrm{aug}})
 &=\boldsymbol f_{\mathrm{aug}}(\theta_{\mathrm{aug}})-\boldsymbol\Phi\,\theta_{\mathrm{aux}},\\
 \boldsymbol\Phi^\top\boldsymbol f^{\perp}_{\mathrm{aug}}(\theta_{\mathrm{aug}})&=0,
\end{aligned}
\label{eq:obc_ours}
\end{equation}
where $\theta_{\mathrm{aux}}\in\R^{10}$ is the coefficient,
$\boldsymbol f^{\perp}_{\mathrm{aug}}:\R^{n_{\mathrm{aug}}}\to\R^{6N}$ the
corrected write, and $\boldsymbol\Phi^{+}$ the Moore-Penrose pseudoinverse, a
left inverse by the full column rank.
Because $\theta_{\mathrm{aux}}$ is differentiated through while
$\boldsymbol\Phi$ is fixed, also
$\boldsymbol\Phi^\top\partial\boldsymbol f^{\perp}_{\mathrm{aug}}/\partial\theta_{\mathrm{aug}}=0$,
so no update of $\theta_{\mathrm{aug}}$ moves the corrected write into
$\ran\boldsymbol\Phi$.
\todo{Missing: a measured check of \eqref{eq:obc_ours} on the thesis data (rank of $\boldsymbol\Phi$ at each rebuild, overlap of the corrected write with $\ran\boldsymbol\Phi$ on held-out tuples). \texttt{DOC/THESIS-RESULTS.md} reports neither, so the overlap display of the earlier draft was removed. Candidate: report both in Section~\ref{sec:results} from the logged validation diagnostics (D-201).}

\subsection{Identification of the corrected model}\label{sec:obc_id}
\ifdraft
\begin{itemize}
\item The expansion point must follow $\xi$; per-epoch refresh adapted from Gy\"or\"ok 2025 Remark~2; reason against every evaluation (cost) and against a fixed point (?) (Gy\"or\"ok 2025 Sec.~4; D-190 item 6; code)
\item Problem display: the corrected physical line in \eqref{eq:aug_loss}, regressor at the simulated point (Gy\"or\"ok 2026 Eq.~(13); code)
\item The added-state line carries no correction (code; D-192)
\item Per phase: Adam, polish with $\xi^\circ$ fixed, S1 and S2 on the corrected model, validation re-solves, prediction with fixed $\xi^\circ$ and $\theta_{\mathrm{aux}}$ (code; D-198)
\item $\mathcal M_{\mathrm{OBC}}$ named; inputs and output as $\mathcal M_{\mathrm U}$ (THESIS-RESULTS step~5; notation (?))
\end{itemize}
\fi

Identifying the corrected model needs an expansion point that follows $\xi$
as the estimate leaves its start.
Gy\"or\"ok et al.\ either update the point at every objective evaluation or fix
it at the nominal parameters~\cite[Sec.~4]{gyorok2025l4dc}.
Updating recomputes the factorisation of the stack at every evaluation.
Fixing it relies on nominal values close to the system.
Their projection is instead recomputed at the start of each epoch from the
current state estimates~\cite[Remark~2]{gyorok2025l4dc}.
We adapt this per-epoch recomputation to the expansion point.
The range of $\boldsymbol\Phi$ turns slowly with $\xi^\circ$, so a more
frequent rebuild changes little.
\todo{Missing: (a) the slow rotation of $\ran\boldsymbol\Phi$ is a measurement on earlier data (D-190 item 6) that no section reports. Candidate: report the principal angle between consecutive bases in Section~\ref{sec:results}. (b) Reason against a fixed point. Candidate: the detuned start is not the nominal value that Gy\"or\"ok et al.\ assume close to the system, so a basis fixed there would protect the directions of the start, not of the estimate (own reading).}

The correction evaluates the regressor at the simulated state and input with
the coefficient of the reference stack, as Gy\"or\"ok et al.\ predict on new
data~\cite[Eq.~(13)]{gyorok2026obc}.
Substituted into the physical state line of \eqref{eq:aug_loss}, it gives the
identification problem
\begin{equation}
\begin{aligned}
 &\min_{\vartheta}\;V(\vartheta)\quad\text{s.t.}\\
 &\tilde x^{\mathrm n}_{k+1|\tau}
  =f^{\mathrm n}_{\mathrm{base}}\bigl(\theta_{\mathrm{base}}(\xi),
   \tilde x^{\mathrm n}_{k|\tau},\hat u^{\mathrm n}_{k|\tau}\bigr)\\
 &\qquad+f^{\mathrm n}_{\mathrm{aug}}\bigl(\theta_{\mathrm{aug}},
   \hat x^{\mathrm n}_{k|\tau},\hat u^{\mathrm n}_{k|\tau}\bigr)
  -\Phi_{\xi^\circ}\bigl(\tilde x^{\mathrm n}_{k|\tau},
   \hat u^{\mathrm n}_{k|\tau}\bigr)\,\theta_{\mathrm{aux}},\\
 &\text{and the other constraints of \eqref{eq:aug_loss}},
\end{aligned}
\label{eq:obc_problem}
\end{equation}
where $V$ and $\vartheta=(\xi,\theta_{\mathrm{aug}},\theta_{\mathrm{enc}})$ are
those of \eqref{eq:aug_loss}, $\theta_{\mathrm{aux}}$ is the coefficient of
\eqref{eq:obc_ours}, solved once per objective evaluation from the current
network and differentiated through, and $\xi^\circ$ is the expansion point,
held fixed within an epoch and set to the current $\xi$ at each epoch
boundary.
The added-state line of \eqref{eq:aug_loss} carries no correction, because the
baseline has no added-state update.

Every phase of Section~\ref{sec:augmentation} identifies
\eqref{eq:obc_problem} in place of \eqref{eq:aug_loss}.
S1 regresses the closed-loop residual of the corrected model and changes only
the head and $\theta^a_{\mathrm{enc}}$, as in Section~\ref{sec:aug_ps2}.
S2 changes $g^{\mathrm n}_{\mathrm{aug}}$, which is uncorrected.
It also moves the shared hidden layers, and so $f^{\mathrm n}_{\mathrm{aug}}$.
The next objective evaluation then solves $\theta_{\mathrm{aux}}$ anew.
S3 and the polish minimise \eqref{eq:obc_problem} as in
Section~\ref{sec:aug_closed_loop}.
Each validation solves $\theta_{\mathrm{aux}}$ from the current network, with
the expansion point of the epoch.
After selection, $\xi^\circ$ is set to the selected estimate and stays there
during the polish.
Prediction holds $\xi^\circ$ and the last solved $\theta_{\mathrm{aux}}$ fixed,
as in~\cite[Eq.~(13)]{gyorok2026obc}.

We call the model of $\mathcal M_{\mathrm{PS2}}$ with the corrected physical
line, identified by \eqref{eq:obc_problem} through S1 to S3,
$\mathcal M_{\mathrm{OBC}}$.
Its inputs and output are those of $\mathcal M_{\mathrm U}$.
\todo{Missing: model notation. Candidate: the calligraphic $\mathcal M_{\mathrm{OBC}}$, since plain $M$ clashes with $M(Y)$ (README open conflict on model names).}

\subsection{Recovery condition and scope}\label{sec:obc_scope}
\ifdraft
\begin{itemize}
\item Interpretable combinations need the unique estimate to be the true one, which depends on the omitted dynamics (Gy\"or\"ok 2026 Sec.~4.1)
\item True one-step discrepancy $\boldsymbol\varepsilon^\star$ on the tuples with a full successor stencil, same normalised frame (D-202, D-203)
\item Proposition~1: on the tuples and to first order, the estimate is $\xi^\star+\boldsymbol\Phi^{+}\boldsymbol\varepsilon^\star$ for every corrected write; recovery iff $\boldsymbol\Phi^\top\boldsymbol\varepsilon^\star=0$ (this work; state-level Gy\"or\"ok Cond.~4, Eq.~(21), without Assumption~6)
\item The condition needs $\theta^\star_{\mathrm{base}}$, so simulation only; input design can meet it (Gy\"or\"ok 2026 Remark~5); whether the thesis data meet it (?)
\item Scope: one-step surrogate of a closed-loop criterion; slice $\bar x=0$ and uncorrected $g^{\mathrm n}_{\mathrm{aug}}$; aggregate and normalised; Theorems 16, 18 do not transfer (D-192; THESIS-RESULTS step~5; Gy\"or\"ok 2026 Thms.~16, 18)
\end{itemize}
\fi

Interpretable combinations also need the unique estimate to be the true one,
which depends on the dynamics the baseline omits.
Let $\xi^\star$ be the training coordinates of the true parameters
$\theta^\star_{\mathrm{base}}$.
On the tuples whose successor sample also admits the difference stencil, the
true one-step discrepancy is
\begin{equation}
\begin{aligned}
 \boldsymbol\varepsilon^\star&=\boldsymbol\chi-\boldsymbol f_{\mathrm{base}}(\xi^\star),\qquad
 \boldsymbol\chi=\col(\chi_i)_{i\in\mathcal I'},\\
 \boldsymbol f_{\mathrm{base}}(\xi)&=\col\bigl(f^{\mathrm n}_{\mathrm{base}}
   (\theta_{\mathrm{base}}(\xi),\tilde x^{\mathrm n}_i,u^{\mathrm n}_i)\bigr)_{i\in\mathcal I'},
\end{aligned}
\label{eq:obc_delta}
\end{equation}
where $\mathcal I'\subset\mathcal I$ drops the last tuple of each record,
$N'=|\mathcal I'|$, $\chi_i\in\R^6$ is $\tilde x^{\mathrm n}_{i+1}$ of
\eqref{eq:obc_stacks}, $\boldsymbol\chi,\boldsymbol\varepsilon^\star\in\R^{6N'}$,
and $\boldsymbol f_{\mathrm{base}}:\R^{10}\to\R^{6N'}$.
The discrepancy collects every part of the recorded successor that the true
baseline does not produce, the error of the velocity estimate included.

\textit{Proposition 1:} Build \eqref{eq:obc_stacks} and \eqref{eq:obc_ours} on
$\mathcal I'$, and let $E(\xi,\theta_{\mathrm{aug}})
=\norm{\boldsymbol\chi-\boldsymbol f_{\mathrm{base}}(\xi)
-\boldsymbol f^{\perp}_{\mathrm{aug}}(\theta_{\mathrm{aug}})}_2^2$ be the
one-step error of the corrected model on these tuples.
Suppose (i) $\boldsymbol\Phi$ has full column rank, and (ii)
$\boldsymbol f_{\mathrm{base}}(\xi)=\boldsymbol f_{\mathrm{base}}(\xi^\star)
+\boldsymbol\Phi(\xi-\xi^\star)$ on a neighbourhood $\mathcal U$ of $\xi^\star$
that contains $\xi^\circ$.
Then every stationary point $(\hat\xi,\hat\theta_{\mathrm{aug}})$ of $E$ with
$\hat\xi\in\mathcal U$ satisfies
$\hat\xi-\xi^\star=\boldsymbol\Phi^{+}\boldsymbol\varepsilon^\star$.
Hence $\hat\xi=\xi^\star$ if and only if
$\boldsymbol\Phi^\top\boldsymbol\varepsilon^\star=0$.

\begin{IEEEproof}
The tuples are data on $\bar x=0$ and $\boldsymbol\Phi$ is fixed, so
$\boldsymbol f^{\perp}_{\mathrm{aug}}$ does not depend on $\xi$.
By (ii), stationarity of $E$ in $\xi$ gives
$\boldsymbol\Phi^\top\bigl(\boldsymbol\varepsilon^\star
-\boldsymbol\Phi(\hat\xi-\xi^\star)
-\boldsymbol f^{\perp}_{\mathrm{aug}}(\hat\theta_{\mathrm{aug}})\bigr)=0$.
The last line of \eqref{eq:obc_ours} removes the learned term, which leaves
$\boldsymbol\Phi^\top\boldsymbol\Phi(\hat\xi-\xi^\star)
=\boldsymbol\Phi^\top\boldsymbol\varepsilon^\star$.
By (i), $\boldsymbol\Phi^\top\boldsymbol\Phi$ is invertible, so
$\hat\xi-\xi^\star=(\boldsymbol\Phi^\top\boldsymbol\Phi)^{-1}
\boldsymbol\Phi^\top\boldsymbol\varepsilon^\star
=\boldsymbol\Phi^{+}\boldsymbol\varepsilon^\star$.
This vanishes exactly when $\boldsymbol\Phi^\top\boldsymbol\varepsilon^\star=0$.
\end{IEEEproof}

Proposition~1 is the state-level counterpart of Condition~4 and Eq.~(21) of
Gy\"or\"ok et al.~\cite[Sec.~4.1]{gyorok2026obc}, for the one-step transition
linearised at $\xi^\circ$.
Unlike their Theorem~7, it needs no exact fit of the data (their
Assumption~6), because the corrected write drops out of the stationarity
condition.
The estimate is therefore the same for every corrected write, which is the
uniqueness Section~\ref{sec:aug_closed_loop} requires, on the tuples and to
first order.
Since \eqref{eq:mdelta_map} is a bijection onto its range,
$\hat\xi=\xi^\star$ means $\hat\theta_{\mathrm{base}}=\theta^\star_{\mathrm{base}}$
whenever $\theta^\star_{\mathrm{base}}$ lies in that range.
The condition needs $\theta^\star_{\mathrm{base}}$, so it can be checked in
simulation only.
When the structure of the omitted term is known, an input can be designed to
meet it~\cite[Remark~5]{gyorok2026obc}.
\todo{Missing: whether the discrepancy of the thesis data meets $\boldsymbol\Phi^\top\boldsymbol\varepsilon^\star=0$. \texttt{DOC/THESIS-RESULTS.md} step~5 states that it does not, but the only measurement (D-190, D-203) used earlier data. Candidate: report $\norm{\boldsymbol\Phi\boldsymbol\Phi^{+}\boldsymbol\varepsilon^\star}/\norm{\boldsymbol\varepsilon^\star}$ on the thesis training records in Section~\ref{sec:results} beside the step~5 table.}

The guarantee is narrower than orthogonality of the simulated output.
Proposition~1 concerns the one-step error on the tuples, whereas
\eqref{eq:obc_problem} minimises a closed-loop multistep output error.
It therefore describes a surrogate, not the bias of $\mathcal M_{\mathrm{OBC}}$.
The tuples lie on $\bar x=0$, so $\theta_{\mathrm{aux}}$ sees the learned write
on that slice only.
In the simulation, the part of the write that acts through a nonzero $\bar x$
has no stacked guarantee.
The uncorrected $g^{\mathrm n}_{\mathrm{aug}}$ also reaches the physical rows
one step later through $f^{\mathrm n}_{\mathrm{aug}}$.
The orthogonality is aggregate over the tuples, not per tuple.
It holds in the normalised coordinates, so the normalisation sets its metric.
The consistency and zero-covariance results of Gy\"or\"ok et
al.~\cite[Thms.~16, 18]{gyorok2026obc} do not transfer.
They assume a stable data-generating
system~\cite[Condition~10]{gyorok2026obc}, which the open-loop gantry is not
(Section~\ref{sec:aug_loop}).
